\section{Graph Validation}

\paragraph{Input and output.}
Graph validation consumes a Graph IR program and either certifies it or reports
a structural error.

\paragraph{Contract.}
The validator checks module inputs and outputs, operand availability,
definition reachability, type compatibility between actual operands and formal
parameters, graph path structure, and backend-relevant invariants such as the
absence of unsupported surface constructs.

\paragraph{Rationale.}
Graph validation is the last generic compiler guard before backend-specific
execution.  It should catch lowering bugs before they become silent numerical
divergences in generated code.

\paragraph{Formal view.}
Graph validation checks a graph program \(G=(M,E)\) against a collection of
well-formedness predicates:
\[
  \mathsf{validateGraph}(G) =
  \mathsf{defs}(G) \land \mathsf{types}(G) \land \mathsf{uses}(G)
  \land \mathsf{paths}(G).
\]
The validator is intentionally stricter than backend code generation: a graph
that fails validation should not be emitted for any backend.

\paragraph{Example.}
If lowering emits a node whose formal parameter expects \code{Tensor[B,S,D]}
but whose actual operand is annotated as \code{Float}, graph validation fails
before torch or tinygrad code is generated.  Similarly, if a node references a
value not produced by module inputs or earlier nodes, validation reports an
undefined operand.
